% ============================================================================
%  COMPANION LaTeX TEMPLATE
%  Follows the strict plain-language + layout rules in companion_style.md.
% ============================================================================

\documentclass[12pt,a4paper]{article}
\usepackage[utf8]{inputenc}
\usepackage[T1]{fontenc}
\usepackage[margin=1in]{geometry}
\usepackage{xcolor}
\usepackage{graphicx}
\usepackage{tcolorbox}
\usepackage{amsmath}
\usepackage{amssymb}
\usepackage{microtype}
\usepackage{booktabs}
\usepackage{adjustbox}
\usepackage{tikz}
\usepackage{enumitem}
\usepackage{fancyhdr}
\usepackage{hyperref}

% --- Colors -----------------------------------------------------------------
\definecolor{navy}{HTML}{06241f}        % Primary brand color (not actually navy but keeping the var name)
\definecolor{cyan}{HTML}{38f0d8}        % Brand highlight
\definecolor{dim}{HTML}{64759c}         % Subtext
\definecolor{intuFrame}{HTML}{38f0d8}   % Intuition box border
\definecolor{workFrame}{HTML}{9b8cff}   % Worked example border
\definecolor{realFrame}{HTML}{ffcb6b}   % Everyday picture border
\definecolor{watchFrame}{HTML}{ff6e81}  % Watch out border
\definecolor{keyFrame}{HTML}{ff6ec7}    % Key takeaway border

% --- Callout boxes ----------------------------------------------------------
\tcbset{
  basebox/.style={
    boxrule=1.5pt,
    arc=6pt,
    boxsep=4pt,
    left=8pt, right=8pt, top=6pt, bottom=6pt,
    colback=white,
    fonttitle=\bfseries\sffamily,
    title filled=false,
    coltitle=navy,
    before skip=16pt,
    after skip=16pt
  }
}

\newtcolorbox{intuition}{
  basebox,
  colframe=intuFrame,
  title=Intuition
}

\newtcolorbox{worked}[1]{
  basebox,
  colframe=workFrame,
  title=Worked Example: #1
}

\newtcolorbox{everyday}{
  basebox,
  colframe=realFrame,
  title=Everyday Picture
}

\newtcolorbox{watchout}{
  basebox,
  colframe=watchFrame,
  title=Watch Out
}

\newtcolorbox{keytake}{
  basebox,
  colframe=keyFrame,
  title=Key Takeaway
}

% --- Environments & helpers -------------------------------------------------
\newenvironment{steps}
  {\begin{enumerate}[label=\textbf{Step \arabic*.}, leftmargin=*]}
  {\end{enumerate}}

\newcommand{\vocab}[1]{\textbf{\textcolor{navy}{#1}}}

\newcommand{\housefig}[2]{
  \begin{figure}[htbp]
    \centering
    \includegraphics[width=0.85\textwidth]{#1}
    \caption{#2}
  \end{figure}
}

% --- Headers & Footers ------------------------------------------------------
\pagestyle{fancy}
\fancyhf{}
\renewcommand{\headrulewidth}{0pt}
\fancyfoot[C]{\thepage}
\fancyhead[L]{\textcolor{dim}{\small NLU \& NLG}}
\fancyhead[R]{\textcolor{dim}{\small Intro to NLP}}

% ============================================================================
%  DOCUMENT START
% ============================================================================
\begin{document}

\begin{center}
  {\LARGE\bfseries\sffamily Natural Language Understanding \& Generation} \\[0.5em]
  {\large A Study Companion} \\[1.5em]
\end{center}

\noindent
\textbf{How to use this guide.}
We build the intuition first, then the math. The \textcolor{realFrame}{\textbf{gold}}
``everyday picture'' box gives you an analogy you already know. The
\textcolor{intuFrame}{\textbf{cyan}} ``intuition'' box tells you the core
mechanism. The \textcolor{workFrame}{\textbf{purple}} ``worked example'' box runs
the numbers. The \textcolor{watchFrame}{\textbf{red}} ``watch out'' box flags the
common mistake. The \textcolor{keyFrame}{\textbf{pink}} ``key takeaway'' box
is the one sentence to remember.

\medskip

% ============================================================================
%  SECTION 1: Natural Language Understanding
% ============================================================================
\section{What is Natural Language Understanding?}

Computers are great at counting numbers. They are terrible at reading books. When a human reads a sentence, they instantly grasp the meaning, the tone, and the context. When a computer reads a sentence, it just sees a long string of letters.

\vocab{Natural Language Understanding} (NLU) is the branch of artificial intelligence that tries to fix this. It is the process of taking raw text and extracting the actual meaning behind it.

\begin{everyday}
Think of NLU like a detective reading a ransom note. The detective does not just read the words; they look for clues. Who wrote it? What do they want? Are they angry? NLU tries to give computers this exact detective skill.
\end{everyday}

\begin{intuition}
To understand text, a computer must break it down into smaller pieces. First it finds the words. Then it looks at how the words connect to each other. Finally, it uses those connections to figure out what the whole sentence means. We call this a pipeline.
\end{intuition}

To help computers act like detectives, we run the text through a series of steps. The first step is usually \vocab{tokenization}. This means splitting a sentence into individual words, or "tokens". The next step might be finding the part of speech for each word (like nouns and verbs).

\begin{worked}{A simple NLU pipeline}
Let's see how a computer might understand the sentence: "Book a flight to Paris."
\begin{steps}
  \item \textbf{Tokenization.} Split the text into tokens: ["Book", "a", "flight", "to", "Paris", "."].
  \item \textbf{Part of Speech.} Tag each word: "Book" (Verb), "a" (Article), "flight" (Noun), "to" (Preposition), "Paris" (Noun).
  \item \textbf{Intent extraction.} The system notices the verb "Book" and the noun "flight". It decides the user's intent is to make a travel reservation.
  \item \textbf{Entity extraction.} The system looks for specific details. It identifies "Paris" as a location entity.
  \item \textbf{Result.} The computer now knows: Intent = BookFlight, Destination = Paris.
\end{steps}
\end{worked}

\housefig{figures/pipeline_flow.png}{The NLU pipeline: raw text goes in, structured meaning comes out.}

\begin{watchout}
Language is messy. A word can mean two different things depending on the context. For example, "I saw a bat." Did you see a flying animal or a piece of wood for hitting baseballs? NLU systems often struggle with this kind of \vocab{ambiguity}.
\end{watchout}

\begin{keytake}
NLU is the process of converting unstructured human language into a structured format that a computer can understand and act upon.
\end{keytake}

\noindent\textbf{ML/AI connection.} Modern NLU systems rarely use manual rules anymore. Instead, they use large neural networks, like Transformers. They predict the meaning directly from the text based on examples seen during training.

% ============================================================================
%  SECTION 2: Natural Language Generation
% ============================================================================
\section{What is Natural Language Generation?}

If NLU is reading, then \vocab{Natural Language Generation} (NLG) is writing. NLG is the process of taking structured data or concepts and turning them into plain, readable human language.

\begin{everyday}
Think of NLG like a news reporter looking at a table of sports scores. The table just has numbers: "Team A: 3, Team B: 1". The reporter writes a story: "Team A won a decisive victory over Team B last night." The reporter turned data into a story.
\end{everyday}

\begin{intuition}
NLG works in reverse from NLU. You start with the meaning you want to convey. Then you select the right words. Finally, you arrange those words into a grammatically correct sentence.
\end{intuition}

Early NLG systems used simple templates, like filling in the blanks. Modern NLG systems use language models. They generate fluent text one word at a time. They predict the next word based on what came before.

\housefig{figures/nlg_process.png}{Language models predict the next word based on context.}

\begin{worked}{Generating a weather report}
Let's turn some simple data into a weather report.
Data: Location = "London", Condition = "Rain", Temp = 15C.
\begin{steps}
  \item \textbf{Content determination.} Decide what to say. We need to mention the location, the rain, and the temperature.
  \item \textbf{Sentence structuring.} Plan the sentence. "It will be [Condition] in [Location] today, with a temperature of [Temp]."
  \item \textbf{Lexicalization.} Choose the specific words. We might change "Rain" to "rainy" to make it sound better.
  \item \textbf{Result.} "It will be rainy in London today, with a temperature of 15 degrees Celsius."
\end{steps}
\end{worked}

\begin{watchout}
Modern language models are very good at generating fluent text, but they can sometimes generate text that is factually incorrect. We call this an \vocab{hallucination}. Just because a computer writes a confident-sounding sentence does not mean it is true.
\end{watchout}

\begin{keytake}
NLG is the process of converting structured data or concepts into fluent, human-readable text. It is the voice of the AI.
\end{keytake}

% ============================================================================
%  GLOSSARY
% ============================================================================
\section*{Glossary}
\addcontentsline{toc}{section}{Glossary}
\noindent\textbf{Key terms.}
\begin{description}[leftmargin=8.5em,style=nextline,font=\normalfont\bfseries\color{navy}]
  \item[NLU] Natural Language Understanding; extracting meaning from text.
  \item[NLG] Natural Language Generation; creating text from data or meaning.
  \item[Tokenization] Splitting text into smaller pieces, like words.
  \item[Intent] What the user is trying to achieve with their sentence.
  \item[Entity] A specific piece of information in the text, like a name or place.
\end{description}

\end{document}
